\chapter{Formal Mathematical Semantics of the EasyCrypt Security Proof Surface}
\label{app:formal-easycrypt-semantics}

\providecommand{\ecfile}[1]{\nolinkurl{#1}}
\providecommand{\ecsym}[1]{\nolinkurl{#1}}
\newcommand{\ECsem}[1]{\llbracket #1 \rrbracket}
\newcommand{\Set}{\mathsf{Set}}
\newcommand{\Dist}{\mathsf{Dist}}
\newcommand{\SubDist}{\mathsf{SubDist}}
\newcommand{\Bool}{\mathsf{Bool}}
\newcommand{\ECState}{\mathsf{State}}
\newcommand{\Prog}{\mathsf{Prog}}
\newcommand{\Kern}{\mathsf{Kern}}

This appendix contains two kinds of notation.  Concrete EasyCrypt identifiers
are written with \ecsym{...} macros and refer to the proof surface discussed in
the main text.  Schematic fragments containing ellipses, such as
\(\tau_1\to\cdots\to\tau_n\) or \ecsym{module M(...)}, are templates for the
EasyCrypt language forms.  They are not quoted theorem headers and should not be
read as exact code correspondence claims.

This appendix fixes the exact mathematical interpretation of the EasyCrypt
objects listed in the generated declaration index.  It is written as a semantic
dictionary: every declaration header in Appendix~\ref{app:easycrypt-declaration-index}
is read by applying one of the clauses below.  The clauses are intentionally
uniform, because the soundness claim of the machine-checked proof is not that
each identifier has a private informal explanation, but that every declaration
is typed, interpreted, and proved in the same formal logic.

\section{Judgmental setting}
\label{app:semantic-judgmental-setting}

Let \(\Gamma\) be the EasyCrypt environment generated by the manifest
\ecfile{provable-security/proof-files.txt}.  The checker establishes judgments
of the form
\[
  \Gamma \vdash d : \mathcal{K},
\]
where \(d\) is a declaration and \(\mathcal{K}\) is its syntactic kind:
\[
  \mathcal{K}\in
  \{\mathsf{Type},\mathsf{Op},\mathsf{Module},\mathsf{Lemma},
    \mathsf{Hoare},\mathsf{Equiv}\}.
\]
The semantic map
\[
  \ECsem{-}_{\Gamma}
\]
sends each checked declaration to a mathematical object.  For the HAETAE
security proof, those objects live in the following universe:
\[
  \Set,\qquad
  \Dist(X),\qquad
  \SubDist(X),\qquad
  X\to Y,\qquad
  X\to\Bool,\qquad
  \Prog,\qquad
  \Kern(X,Y),
\]
where \(\Dist(X)\) denotes probability distributions over \(X\),
\(\SubDist(X)\) denotes subdistributions over \(X\), and
\(\Kern(X,Y)=X\to\SubDist(Y)\) denotes a general probabilistic program kernel.
Losslessness lemmas identify the special cases whose kernels have total mass one
and can therefore be read as \(X\to\Dist(Y)\).

\section{Types}
\label{app:semantic-types}

The exact EasyCrypt forms used by the manifest are:

\begin{lstlisting}[language=EasyCrypt]
type t.
type t = tau.
type t = [ C1 | ... | Cn ].
type t = tau1 * ... * taun.
\end{lstlisting}

The mathematical reading is:
\[
  \Gamma \vdash \mathtt{type}\ t : \mathsf{Type}
  \quad\Longmapsto\quad
  \ECsem{t}_{\Gamma}=X_t\in\Set.
\]
If \(t\) is abstract, \(X_t\) is an uninterpreted carrier set.  If \(t\) is a
finite variant type, then
\[
  \ECsem{[C_1|\cdots|C_n]}_{\Gamma}
  =
  \{C_1,\ldots,C_n\}.
\]
If \(t\) is a product type, then
\[
  \ECsem{\tau_1*\cdots*\tau_n}_{\Gamma}
  =
  \ECsem{\tau_1}_{\Gamma}\times\cdots\times\ECsem{\tau_n}_{\Gamma}.
\]
If \(t\) is a list type, then
\[
  \ECsem{\tau\ \mathsf{list}}_{\Gamma}
  =
  \ECsem{\tau}_{\Gamma}^{*}.
\]

For the HAETAE proof this gives the following exact mathematical classes:
\[
\begin{array}{rcl}
\ECsem{\mathtt{pkey}} &=& \mathsf{PK},\\
\ECsem{\mathtt{skey}} &=& \mathsf{SK},\\
\ECsem{\mathtt{message}} &=& \mathsf{Msg},\\
\ECsem{\mathtt{context}} &=& \mathsf{Ctx},\\
\ECsem{\mathtt{signature}} &=& \mathsf{Sig},\\
\ECsem{\mathtt{ro\_query}} &=& \mathsf{Q}_{\mathsf{RO}},\\
\ECsem{\mathtt{ro\_output}} &=& \mathsf{Y}_{\mathsf{RO}},\\
\ECsem{\mathtt{mode}} &=& \{\mathsf{Mode2},\mathsf{Mode3},\mathsf{Mode5}\},\\
\ECsem{\mathtt{poly}} &=& \mathbb{Z}^{*}\ \text{with well-formed subset } \mathbb{Z}^{256},\\
\ECsem{\mathtt{matrix}} &=& (\mathbb{Z}^{*})^{* *}\ \text{with mode-indexed well-formed subset}.
\end{array}
\]
The proof scripts use predicates such as \ecsym{poly\_wf},
\ecsym{polyveck\_wf}, and \ecsym{challenge\_wf} to restrict list carriers to
the intended finite-dimensional algebraic domains.

\section{Operations}
\label{app:semantic-ops}

The exact EasyCrypt forms used by the manifest are:

\begin{lstlisting}[language=EasyCrypt]
op c : tau.
op f (x1 : tau1) ... (xn : taun) : tau = body.
op d : tau distr.
op [lossless] d : tau distr.
abbrev a = body.
\end{lstlisting}

The mathematical reading is:
\[
  \Gamma\vdash \mathtt{op}\ f : \tau_1 \to \cdots \to \tau_n \to \tau
  :\mathsf{Op}
  \quad\Longmapsto\quad
  \ECsem{f}_{\Gamma}:
  \ECsem{\tau_1}_{\Gamma}\times\cdots\times\ECsem{\tau_n}_{\Gamma}
  \to \ECsem{\tau}_{\Gamma}.
\]
The codomain determines the mathematical role:
\[
\begin{array}{rcll}
\tau=\mathsf{bool}
&\Longrightarrow&
\ECsem{f}\subseteq
\ECsem{\tau_1}\times\cdots\times\ECsem{\tau_n},
& \text{predicate or event},\\[0.3em]
\tau=\sigma\ \mathsf{distr}
&\Longrightarrow&
\ECsem{f}:
\ECsem{\tau_1}\times\cdots\times\ECsem{\tau_n}
\to \SubDist(\ECsem{\sigma}),
& \text{sampler or law; losslessness gives a distribution},\\[0.3em]
\tau=\mathsf{real}
&\Longrightarrow&
\ECsem{f}\in\mathbb{R}
\text{ or }
\ECsem{f}:\cdots\to\mathbb{R},
& \text{advantage or loss},\\[0.3em]
\tau=\sigma_1*\cdots*\sigma_m
&\Longrightarrow&
\ECsem{f}:\cdots\to
\ECsem{\sigma_1}\times\cdots\times\ECsem{\sigma_m},
& \text{constructor}.
\end{array}
\]

The following examples are exact EasyCrypt declarations from the proof surface
and their mathematical readings.

\begin{lstlisting}[language=EasyCrypt]
op fresh_msg (qs : query list) (m : message) (ctx : context) : bool =
  ! mem qs (m, ctx).
\end{lstlisting}
\[
  \ECsem{\mathtt{fresh\_msg}}(Q,m,c)
  \equiv
  (m,c)\notin Q.
\]

\begin{lstlisting}[language=EasyCrypt]
op sampler_expand_query (coins : random_coins) : ro_query =
  SamplerExpandQuery coins.
\end{lstlisting}
\[
  \ECsem{\mathtt{sampler\_expand\_query}}:
  \mathsf{Coins}\to\mathsf{Q}_{\mathsf{RO}},
  \qquad
  r\mapsto \mathsf{SamplerExpandQuery}(r).
\]

\begin{lstlisting}[language=EasyCrypt]
op counted_rom_programming_loss_term =
  counted_rom_collision_term + counted_rom_prequery_term
  + counted_rom_min_entropy_term.
\end{lstlisting}
\[
  L_{\mathsf{ROM}}
  =
  L_{\mathsf{coll}}
  +
  L_{\mathsf{pre}}
  +
  L_{\mathsf{me}}.
\]

\begin{lstlisting}[language=EasyCrypt]
op dsigning_randomness : random_coins distr =
  dmap dsigning_randomness_source signing_randomness_from_source.
\end{lstlisting}
\[
  \mu_{\mathsf{coins}}
  =
  (\mathsf{signing\_randomness\_from\_source})_{\#}
  \mu_{\mathsf{source}},
\]
where \((-)_{\#}\) denotes the push-forward of a distribution.

\section{Modules and module types}
\label{app:semantic-modules}

The exact EasyCrypt forms used by the manifest are:

\begin{lstlisting}[language=EasyCrypt]
module type I = { proc p(x : tau) : sigma }.
module M(A : I1, O : I2) = { ... }.
clone import T as U with type t <- tau, op f <- g.
\end{lstlisting}

A module type is an interface:
\[
  \ECsem{\mathtt{module\ type}\ I}_{\Gamma}
  =
  \mathcal{I}
  =
  \{(p_i:X_i\to\SubDist(Y_i))_i\}.
\]
A concrete module implementing this interface is a tuple of probabilistic
procedures:
\[
  \ECsem{M}_{\Gamma}\in\mathcal{I}.
\]
A functorial module is a higher-order construction:
\[
  \ECsem{M(A,O)}_{\Gamma}
  :
  \ECsem{A}\times\ECsem{O}
  \to
  \Prog.
\]
For a security game module \(G(A)\), the procedure \ecsym{main} denotes a
Bernoulli experiment:
\[
  \ECsem{G(A).\mathtt{main}}:
  \mathbf{1}\to\SubDist(\Bool),
  \qquad
  \Pr[G(A).\mathtt{main}()=\mathsf{true}]
  =
  \Adv_G(A).
\]

For HAETAE, the key module interpretations are:
\[
\begin{array}{rcl}
\ECsem{\mathtt{EUF\_CMA}(H,S,A)} &=&
\text{chosen-message unforgeability experiment},\\
\ECsem{\mathtt{UF\_NMA}(H,S,A)} &=&
\text{no-message unforgeability experiment},\\
\ECsem{\mathtt{CountedLazyROM}} &=&
\text{lazy random oracle with query and bad-event logs},\\
\ECsem{\mathtt{ROSigningAttemptPaperSimSampler}} &=&
\text{signing-attempt sampler kernel},\\
\ECsem{\mathtt{ROExactHyperballPaperSimSampler}} &=&
\text{exact-hyperball sampler kernel},\\
\ECsem{\mathtt{BimodalMSIS\_As\_ModuleSIS}} &=&
\text{reduction from bimodal self-target MSIS to Module-SIS}.
\end{array}
\]

\section{Lemmas}
\label{app:semantic-lemmas}

The exact EasyCrypt form is:

\begin{lstlisting}[language=EasyCrypt]
lemma L (x1 : tau1) ... (xn : taun) : Phi.
proof.
  ...
qed.
\end{lstlisting}

The mathematical reading is the theorem judgment
\[
  \Gamma\vdash L:\Phi
  \quad\Longmapsto\quad
  \Gamma\models
  \forall x_1\in\ECsem{\tau_1},\ldots,x_n\in\ECsem{\tau_n}.
  \ECsem{\Phi}.
\]
The proof script between \ecsym{proof.} and \ecsym{qed.} is a kernel-checked
derivation of this proposition.

The main lemma families have the following exact symbolic forms.

\paragraph{Well-formedness lemmas.}
An EasyCrypt lemma such as
\begin{lstlisting}[language=EasyCrypt]
lemma poly_add_wf p r : poly_wf (poly_add p r).
\end{lstlisting}
means
\[
  \forall p,r\in\mathsf{Poly}.
  \mathsf{WF}_{\mathsf{poly}}(p+r).
\]
When the definition of \ecsym{poly\_wf} is unfolded, this is a dimension
preservation theorem for the list representation of polynomials.

\paragraph{Losslessness lemmas.}
An EasyCrypt lemma such as
\begin{lstlisting}[language=EasyCrypt]
lemma signing_coin_distribution_lossless : is_lossless drandom_coins.
\end{lstlisting}
means
\[
  \sum_{r\in\mathsf{Coins}}\mu_{\mathsf{coins}}(r)=1.
\]

\paragraph{Game-hop lemmas.}
A game-hop lemma has mathematical form
\[
  \Pr[G_i^A\Rightarrow 1]
  \le
  \Pr[G_{i+1}^{B(A)}\Rightarrow 1]+L_i.
\]
The EasyCrypt statement names the concrete modules and loss term; the proof
shows that the displayed probability inequality follows from an equivalence,
Hoare invariant, or bad-event bound.

\paragraph{Reduction lemmas.}
A reduction lemma has mathematical form
\[
  \Adv_{\mathsf{HAETAE}}^{G_i}(A)
  \le
  \Adv_{\mathsf{Hard}}(B_A)+L.
\]
Here \(B_A\) is either an explicit EasyCrypt adapter module or an abstract
solver/reduction interface.  For an explicit adapter lemma, the proof connects
source-game success to target-game success up to loss \(L\).  For an abstract
interface lemma, that connection is an explicit antecedent supplied to the
composition theorem.

\paragraph{Top-level lemmas.}
The final checked implication family in \ecfile{HAETAE\_Security.ec} has the
fixed-ledger mathematical shape
\[
  \mathcal{H}_{\mathsf{hop}}
  \wedge
  \mathcal{H}_{\mathsf{hard}}
  \Longrightarrow
  \Adv_{\mathsf{ECModel},\rom}^{\eufcma}(A)
  \le
  B_{\mathsf{counted\ ROM}}
  +
  L_{\mathsf{ModuleSIS}}
  +
  L_{\mathsf{MLWE}}
  +
  L_{\mathsf{bimodal}\to\mathsf{MSIS}},
\]
where \(B_{\mathsf{counted\ ROM}}\) is the sum of the checked collision,
prequery, and min-entropy terms in \ecfile{HAETAE\_Reductions.ec}, and
\(L_{\mathsf{ModuleSIS}}\) and \(L_{\mathsf{MLWE}}\) are the declared hardness
terms obtained from the hardness premises.  The separate
paper-style rejection/Fiat--Shamir-with-aborts expression
\ecsym{haetae\_euf\_bound} is not an additional summand in the counted-ROM
theorem \ecsym{haetae\_euf\_counted\_rom\_bound}.  The antecedents
\(\mathcal{H}_{\mathsf{hop}}\) and \(\mathcal{H}_{\mathsf{hard}}\) are not
implicit: they are the explicit hop/reduction and hardness premises of the
EasyCrypt lemma being cited.

\section{Hoare, probabilistic Hoare, and equivalence judgments}
\label{app:semantic-hoare-equiv}

The exact EasyCrypt proof forms used by the manifest are:

\begin{lstlisting}[language=EasyCrypt]
hoare.
phoare.
equiv.
byequiv (: P ==> Q) => //.
\end{lstlisting}

The mathematical judgments are:
\[
  \vdash\{P\}\ c\ \{Q\}
  \quad\text{for ordinary Hoare proofs,}
\]
\[
  \vdash\Pr[c(s):Q]\ \bowtie\ p
  \quad(\bowtie\in\{=,\le,<,\ge,>\})
  \quad\text{for probabilistic Hoare proofs,}
\]
and
\[
  \vdash c_1 \sim c_2 :
  R_{\mathsf{pre}}\Rightarrow R_{\mathsf{post}}
  \quad\text{for relational equivalence proofs.}
\]

In the HAETAE game hops, these judgments are used in three mathematically
distinct ways.

\paragraph{Invariant preservation.}
For a signing oracle \(O.\mathtt{sign}\), a Hoare proof establishes
\[
  \{I(s)\}\ O.\mathtt{sign}(m,c)\ \{I(s')\wedge Q(s,s',m,c)\}.
\]
Here \(I\) may be transcript-log consistency, active public-key agreement,
budget discipline, or sampler-ROM coverage.

\paragraph{Coupling of adjacent games.}
An equivalence proof establishes a relation between two runs:
\[
  G_i^A \sim G_{i+1}^A :
  R_0 \Rightarrow
  \bigl(res_i=res_{i+1}\bigr)\vee\mathsf{bad}.
\]
This is the formal game-playing statement that the two experiments are identical
until bad.

\paragraph{Fundamental lemma of game playing.}
If the equivalence gives equality outside \(\mathsf{bad}\), the probability
difference is bounded by
\[
  \left|
  \Pr[G_i^A\Rightarrow 1]
  -
  \Pr[G_{i+1}^A\Rightarrow 1]
  \right|
  \le
  \Pr[G_i^A:\mathsf{bad}].
\]
In the HAETAE sampler lifting proof, the relevant bad event is
\[
  \mathsf{Bad}_{\mathsf{prequery}}
  =
  \{\mathsf{sampler\_expand\_query}(R)
    \in Q_A\cup Q_{\mathsf{sampler}}\},
\]
where \(R\leftarrow\mu_{\mathsf{coins}}\) is fresh signing randomness,
\(Q_A\) is the adversary hash-query component of the log, and
\(Q_{\mathsf{sampler}}\) is the sampler self-log from earlier signing calls.
For one fresh signing call, the corresponding bound uses the combined
adversary-hash and sampler-self-query budget.  The checked game-level NMA
statement also applies a union bound over signing calls:
\[
  \Pr[\mathsf{Bad}_{\mathsf{prequery}}]
  \le
  q_S\cdot\frac{q_H+q_S}{|\mathcal C|}.
\]
In EasyCrypt, the checked NMA sampler-prequery fundamental-lemma bound has
right-hand side
\[
  \ec{rom\_signature\_query\_budget}
  \cdot
  \frac{
    \ec{rom\_hash\_query\_budget}
    +
    \ec{rom\_signature\_query\_budget}}
  {\ec{challenge\_support\_cardinality\_lower\_bound}} .
\]

\section{Exact correspondence of proof-object classes}
\label{app:proof-object-exact-correspondence}

\begin{longtable}{p{0.18\textwidth}p{0.34\textwidth}p{0.38\textwidth}}
\caption{Exact mathematical correspondence for proof-object classes.}
\label{tab:proof-object-exact-correspondence}\\
\toprule
Class & EasyCrypt code form & Mathematical object \\
\midrule
\endfirsthead
\toprule
Class & EasyCrypt code form & Mathematical object \\
\midrule
\endhead
Types &
\ecsym{type t.} &
Carrier set \(X_t\in\Set\). \\
Ops &
\ecsym{op f : tau1 -> tau2.} &
Total function \(f:\ECsem{\tau_1}\to\ECsem{\tau_2}\). \\
Predicate ops &
\ecsym{op P : tau -> bool.} &
Subset/event \(P\subseteq\ECsem{\tau}\). \\
Distribution ops &
\ecsym{op d : tau distr.} &
Subprobability law \(\mu_d\in\SubDist(\ECsem{\tau})\); a
\ecsym{[lossless]} annotation or \ecsym{is\_lossless} lemma upgrades it to
\(\Dist(\ECsem{\tau})\). \\
Loss ops &
\ecsym{op L : real.} &
Real bound term \(L\in\mathbb{R}_{\ge0}\), with non-negativity proved separately. \\
Module types &
\ecsym{module type I = \{...\}.} &
Interface \(\mathcal I\) of admissible probabilistic algorithms. \\
Modules &
\ecsym{module M(...) = \{...\}.} &
Probabilistic algorithm, oracle, game, sampler, or reduction. \\
Game modules &
\ecsym{module G(A) = \{ proc main ... \}.} &
Experiment \(G^A\) with success probability \(\Pr[G^A\Rightarrow1]\). \\
Lemmas &
\ecsym{lemma L : Phi.} &
Theorem \(\Gamma\models\Phi\). \\
Hoare &
\ecsym{hoare.} &
Program theorem \(\vdash\{P\}c\{Q\}\). \\
Probabilistic Hoare &
\ecsym{phoare.} &
Probability theorem \(\vdash\Pr[c:Q]\bowtie p\). \\
Equivalence &
\ecsym{equiv.} or \ecsym{byequiv ...} &
Relational/coupling theorem \(c_1\sim c_2:R\Rightarrow S\). \\
\bottomrule
\end{longtable}

Appendix~\ref{app:easycrypt-declaration-index} gives the extracted declaration
headers to which these clauses apply.  Therefore each listed HAETAE proof
object has both a concrete EasyCrypt declaration and a precise mathematical
interpretation.
